\subsection{Composition of Functions} 
As we develop a greater understanding of functions, we will want to develop methods of combining functions. The first of these is \term{composition}.
\begin{definition}[Composition of Functions]
The \term{composition} of two functions, written $$f\circ g$$ and read ``$f$ composed with $g$'' is the function such that, for all $x$ for which it is defined, $(f\circ g)(x)=f(g(x))$.
\end{definition}
\begin{note} We use the symbol $\circ$ to denote composition, and the symbol $\cdot$ to denote multiplication. Make sure you can tell the difference between these symbols in your handwriting.\end{note}
\begin{example}
Let $f=\Set{(1,4),(2,1),(3,3),(4,2)}$ and $g=\Set{(1,3),(2,1),(3,4),(4,2)}$. Find each of the following.
\begin{lpicvsplit}{3}{5}
\foreach \x/\lbl in {0/(f\circ g)(2),1/(f\circ g)(2),2/(g\circ f)(3)}
	\regtextleft{\x*\value{fcols}+\x+1}{-1}{$\lbl={}$};
\end{lpicvsplit}
\end{example}
\begin{example}
For the functions $f$ and $g$, graphed below, find each of the following.
\begin{center}
\begin{tikzpicture}[x=0.5\linespace,y=0.5\linespace]
\useasboundingbox (-5.5,-5.5) rectangle (5.5,5.5);
\setgraphsquare{5}
\GraphSmall
\draw[graph,red] ({2-sqrt(13.8)},-5) parabola bend (2,29/8) (5,-2);
\draw[graphend,red] (2,29/8) node[anchor=south west,font=\small,inner sep=0.5pt]{$y=f(x)$};
\foreach \x/\y in {-1/-2,1/3,3/3} \draw[dotfull,red] (\x,\y) circle (\dotsize);
\linepoints(-2,3)(1,1)
\draw[graphend] (-2,3) node[anchor=north east,font=\small,inner sep=0.5pt]{$y=g(x)$};
\foreach \x/\y in {-2/3,1/1,4/-1} \draw[dotfull] (\x,\y) circle (\dotsize);
\end{tikzpicture}
\end{center}
\begin{lpicvsplit}{3}{5}
\foreach \x/\lbl in {0/(f\circ g)(1),1/(f\circ g)(4),2/(g\circ f)(-1)}
	\regtextleft{\x*\value{fcols}+\x+1}{-1}{$\lbl={}$};
\end{lpicvsplit}
\end{example}